\begin{abstract}
Every sequence of generalized-canonical flips on compact K\"ahler
fourfolds with rational boundary and fixed rational analytically nef
b-divisor terminates when the small morphisms are projective. We prove
this for normal globally Weil $\Q$-factorial models with boundary
coefficients less than one, allowing exceptional log discrepancies equal
to one. No scaling rule or pseudo-effectivity assumption is required.
The theorem applies to existing projective small diagrams with the
specified opposite ample signs.
\end{abstract}

\section{Introduction and the main theorem}\label{sec:introduction}

A flip replaces a small contraction on which an adjoint divisor is
negative by a small contraction on which its transform is positive.
Termination asks whether this local improvement can continue
indefinitely. Generalized pairs enlarge the adjoint from $K_X+B$ to
$K_X+B+M_X$, where $M_X$ is the trace of nef data on a higher
birational model. The nef contribution can change discrepancies even
where the ordinary boundary is absent. On a compact K\"ahler space,
one must also distinguish analytic nefness from nonnegative degrees
on curves.

We prove termination in the generalized-canonical case in dimension
four, for rational nef data fixed throughout the sequence. The proof
does not choose the contractions. It applies to every sequence of
projective small diagrams with the required opposite ample signs.
Exceptional log discrepancies may equal one. The resulting canonical
singularities need not be smooth in codimension two, so the argument
must establish smoothness precisely along the flipped surfaces that
it counts. We first specify the global analytic category of the theorem.

\subsection{The data and the statement}

All spaces are complex, and a fourfold is irreducible and
four-dimensional. A K\"ahler form on a normal space is
given by smooth strictly plurisubharmonic local potentials in analytic
embeddings. A class in the real Bott--Chern space
$H^{1,1}_{\mathrm{BC}}(X,\R)$ is \emph{analytically nef} if it lies
in the closure of the K\"ahler cone. For a rational line bundle or
a $\Q$-Cartier divisor, this means that its first Chern class is nef
in this sense.

We call a normal compact space \emph{globally Weil $\Q$-factorial} if
every prime Weil divisor defined on the whole space is $\Q$-Cartier
and some positive reflexive power of its canonical sheaf is a line
bundle. This condition does not assert factoriality of the local
analytic rings, or $\Q$-Cartierness of every rank-one reflexive sheaf.
Compatible actual canonical divisor representatives on the birational
models are part of the supplied data. We do not infer their existence
from the canonical rational line-bundle condition alone.

Fix a projective birational morphism
\[
 \pi:X'\longrightarrow X
\]
of normal compact K\"ahler spaces and an analytically nef
$\Q$-Cartier divisor $M'$ on $X'$. These data define a rational
b-divisor $\bM$: on a model dominating $X'$ its trace is the pullback
of $M'$, and on another model its trace is the pushforward from a
common higher model. We write $M_T$ for the trace on a model $T$.
The divisor $M'$ need not be effective. Replacing $X'$ by a higher
model and $M'$ by its pullback leaves $\bM$ unchanged.

Let $B\geq0$ be a rational divisor of finite support with coefficients
less than one, and assume $D=K_X+B+M_X$ is $\Q$-Cartier. On a smooth
model $p:W\to X$ dominating $X'$, define the actual rational divisor
$\Delta_W$ by
\begin{equation}\label{eq:structure}
 K_W+\Delta_W+M_W=p^*(K_X+B+M_X).
\end{equation}
For a prime divisor $E$ on $W$, its \emph{generalized log discrepancy}
is
\[
 a(E;X,B+\bM)=1-\operatorname{coeff}_E\Delta_W.
\]
Here a \emph{prime over $X$}, or a \emph{divisorial place}, is a prime
divisor on a proper bimeromorphic model, identified with its strict
transforms on common higher models. Its centre is its reduced
irreducible image, and it is \emph{exceptional} if this image has
codimension at least two. Equation~\eqref{eq:structure} on higher
models defines its discrepancy independently of the carrying model.
The pair is \emph{generalized klt} if all these discrepancies are
positive, and \emph{generalized canonical} if the discrepancies of
exceptional primes are at least one. Replacing the latter inequality
by strict inequality gives \emph{generalized terminal}. Throughout
the theorem we impose positivity for every prime as well as the
canonical bound for exceptional primes. Accordingly, exceptional
coefficients of $\Delta_W$ may be zero or negative.

We use \emph{projective} for a proper morphism admitting a relatively ample
holomorphic line bundle. A $\Q$-Cartier divisor is relatively ample
if a positive Cartier multiple is so. A proper bimeromorphic
morphism is \emph{small} if its exceptional locus contains no prime
divisor. A small birational map identifies the prime divisors on its
two models, and hence defines strict transforms of boundaries and
canonical divisors.

\begin{theorem}\label{thm:termination}
Let $X_0$ be a normal irreducible globally Weil $\Q$-factorial compact
K\"ahler fourfold, and let $B_0\geq0$ be a rational divisor with
coefficients in $[0,1)$. Fix compatible actual canonical divisors on
the birational models. Let $\bM$ be the rational b-divisor represented
by an analytically nef $\Q$-Cartier divisor $M'$ on a projective
birational morphism $X'\to X_0$ of normal compact K\"ahler spaces.
Assume
\[
 \begin{aligned}
 a(E;X_0,B_0+\bM)&>0&&\text{for every prime over $X_0$},\\
 a(E;X_0,B_0+\bM)&\geq1&&\text{for every exceptional prime}.
 \end{aligned}
\]
Consider a sequence of normal globally Weil $\Q$-factorial compact
K\"ahler fourfolds
\[
 X_0\dashrightarrow X_1\dashrightarrow X_2\dashrightarrow\cdots .
\]
On $X_i$, let $B_i$ be the strict transform of $B_0$, let $M_i$ be
the trace of this same $\bM$, and suppose
$D_i=K_{X_i}+B_i+M_i$ is $\Q$-Cartier. Suppose that every step is
given by a diagram
\begin{equation}\label{eq:step-diagram}
 X_i\xrightarrow{\ f_i\ }Z_i
       \xleftarrow{\ f_i^+\ }X_{i+1}
\end{equation}
of projective small bimeromorphic morphisms with connected fibres,
where $Z_i$ is normal compact K\"ahler, both morphisms are
nonisomorphisms, $-D_i$ is $f_i$-ample, and $D_{i+1}$ is
$f_i^+$-ample. Then the sequence is finite.
\end{theorem}

The canonical discrepancy bound and positivity are assumed only on
the initial pair; both persist by the comparison proved below.
The theorem imposes no pseudo-effectivity or scaling condition and
no relative Picard-number restriction. It concerns finiteness of
given diagrams; existence of a desired next step is a separate question.

\subsection{Preceding results and the present argument}

In the relative projective algebraic setting, Kawamata, Matsuda,
and Matsuki proved termination for ordinary terminal
$\Q$-factorial fourfolds with empty boundary. Their fourfold argument
combines Shokurov's discrepancy difficulty with a descent of surface
classes \cite[Theorem~5-1-15 and its proof]{KMM1987}.
Fujino extended the discrepancy and surface-count method to ordinary
canonical projective fourfold pairs with effective rational boundary
\cite{Fujino2004}. The corrected statement with
$\lfloor B\rfloor=0$ is \cite[Theorem~5.1]{Fujino2005Addendum}; its
flipping contractions need not have relative Picard number one.
Two parts of this work directly underlie our proof. First, the target
of a canonical flip is terminal near general points of its
codimension-two flipped locus, which permits the blowup discrepancy
calculation \cite[Lemma~2.2]{Fujino2004}. Second, the corrected
low-discrepancy count makes its exclusions by centre, so that they
apply to repeated blowups above a boundary surface and not only to
its first blowup
\cite[Proposition~3.1 and Lemma~3.2]{Fujino2005Addendum}.

Generalized pairs retain nef information on a higher birational
model; see the generalized polarized pairs of
Birkar and Zhang \cite[Definition~1.4 and Section~4]{BirkarZhang2016}.
Chen and Tsakanikas proved termination of every sequence of flips for
pseudo-effective four-dimensional NQC log canonical generalized pairs
in the relative projective algebraic setting
\cite[Theorem~1.1]{ChenTsakanikas2023}; see also
\cite[Theorem~4.4]{Moraga2025}.
Here NQC means that the nef data are finite nonnegative real
combinations of nef $\Q$-Cartier data.
The rational nef datum used here is NQC in the algebraic setting.
Their result allows log canonical singularities and real NQC data.
Theorem~\ref{thm:termination} instead concerns canonical singularities
and fixed rational data, permits nonprojective compact K\"ahler spaces,
and applies without pseudo-effectivity. The smooth-centre discrepancy
formula for generalized pairs also appears in
\cite[Lemma~2.9]{ChenTsakanikas2023}.

In the K\"ahler setting, Das, Hacon, and P\u{a}un prove termination
for ordinary dlt compact K\"ahler fourfold pairs whose adjoint is rationally linearly
equivalent to an effective divisor, assuming the successive models
remain K\"ahler \cite[Theorem~6.6]{DHP2024}.
Hacon and Xie prove termination of a minimal model program with
scaling for strongly $\Q$-factorial compact K\"ahler generalized klt
pairs whose boundary plus nef part is big, in arbitrary dimension
\cite[Theorem~1.5]{HX2026}. The present theorem treats every given
sequence in its canonical rational setting, without either of these
positivity or scaling conditions.

The ordinary terminal K\"ahler fourfold argument in
\cite[Section~4]{OrdinaryKahler} uses spans of analytic cycle classes
in place of algebraic cycle classes. We retain this organization and
establish the generalized-canonical version. Its essential points are
local ordinary terminality along flipped surfaces, a finite count of
places below log discrepancy two allowing zero exceptional boundary
coefficients, and a fixed denominator for witness discrepancies at
smooth general points. The cycle comparison is proved for compact
reduced analytic spaces, so it applies directly to boundary components
even when they are nonnormal. The proof reproduces these arguments
in full and invokes no termination theorem. Its analytic negativity
input is Wang's lemma \cite[Lemma~1.3 and Appendices~A--B]{Wang2021}.
The references to \cite{OrdinaryKahler} concern only this negativity
statement, terminal smoothness in codimension two, and the cycle
comparison. No canonical bundle formula or absolute supporting-contraction
theorem from that paper is used: the small projective diagrams are part
of the hypotheses of Theorem~\ref{thm:termination}.

\subsection{Proof strategy}

Dimension four links the two parts of the proof. Smallness bounds both
exceptional loci by dimension two, and the opposite ample signs force
their dimensions to sum to at least three. Once the target locus
contains no surface, the source locus must therefore contain one.
A surface also has codimension two: its blowup at a smooth general
point has ordinary log discrepancy two. This connects the discrepancy
threshold to the surface classes that will be counted.

Surfaces contained in the source and target exceptional loci are
called \emph{flipping} and \emph{flipped surfaces}, respectively.
First we prove strict discrepancy increase for every place centred
in either exceptional locus. Thus a place exceptional over the target
and centred in a flipped surface has log discrepancy strictly greater
than one, even when its source discrepancy equals one. Effectivity of
the boundary and of the difference between the pulled-back nef trace
and the nef datum then give ordinary terminality near a general point
of that surface. Terminal smoothness in codimension two makes the
target smooth there. This is the precise smoothness needed in the proof.

Blowing up a flipped surface at such points supplies a global prime
divisor, called its \emph{witness}. Choose one integer $m>0$ such that
$mM'$ is Cartier and $mB_0$ has integral coefficients. The integral
Weil divisor $mM_T$ is Cartier near the smooth general point of every
flipped surface on a target $T$. The witness calculation consequently
puts its target log discrepancy in
\[
 (1,2]\cap m^{-1}\Z.
\]
Its discrepancy strictly increases at each use, so each fixed place
can serve as a witness at most $m$ times. This count uses no uniform
bound on the Cartier indices of the nef traces elsewhere on the models.

We combine this finite set of values with a finite set of possible
places, treating the positive boundary coefficients in decreasing
order. If a flipped surface lies in a component of coefficient $b$,
its witness has source discrepancy strictly below $2-b$.
On the first model of a fixed tail, monotonicity keeps its discrepancy
below $2-b$. The low-place lemma puts these divisors in a finite set,
except possibly those centred on surfaces in a positive boundary
component of coefficient $d$, with discrepancy at least $2-d$. Thus $d>b$.
The induction has already
excluded flipping surfaces in those higher-coefficient components,
so these centres persist through the common isomorphic opens and
cannot become flipped surfaces. At each coefficient stage the
possible divisorial witnesses thus belong to a fixed finite set.

Once these witnesses are exhausted, the target exceptional locus
contains no surface in the boundary components of the current
coefficient. On each such component,
Borel--Moore localization compares
its analytic surface classes with those of the source component.
Every removed source surface gives a nonzero kernel class by
K\"ahler positivity. The rank consequently drops at each remaining
step affecting such a surface. After the boundary components have
been treated, the witness argument with threshold two excludes every
remaining flipped surface. The same cycle comparison on the fourfolds
then gives strict rank loss at every step.

Section~\ref{sec:foundations} fixes the analytic foundations and
common models. Section~\ref{sec:comparison} proves the strict
one-step comparison. Section~\ref{sec:low-places} establishes the
low-place count. Section~\ref{sec:witnesses} proves generic
terminality along flipped surfaces and constructs witnesses with
finitely many possible discrepancy values. Section~\ref{sec:cycles} proves the
cycle-rank comparison, and Section~\ref{sec:termination} combines
these ingredients.
